\notepage{80}{Scattering amplitude and differential cross section}
Keeping only the outgoing radial wave gives
\begin{align*}
 \psi_{\rm sc}&\sim\frac{\ee^{\ii kr}}{kr}
 \sum_{\ell,m}(-\ii)^\ell(-B_{\ell m})Y_\ell^m(\theta,\phi)\\
 &=\frac{\ee^{\ii kr}}r f(\theta,\phi),
\end{align*}
which defines the scattering amplitude $f$. Thus
\[
 \boxed{\psi_k\sim\ee^{\ii kz}
              +f(\theta,\phi)\frac{\ee^{\ii kr}}r.}
\]
For a finite beam of transverse extent $\rho_{\max}$, the notes distinguish
the forward incident-beam cone, $\theta\lesssim\rho_{\max}/r$, from
detectors outside that cone. The incident plane-wave current is
\begin{align*}
 \mathbf j_{\rm inc}
 &=\frac{\hbar}{2\mu\ii}
 \left(\ee^{-\ii kz}\boldsymbol\nabla\ee^{\ii kz}
      -\ee^{\ii kz}\boldsymbol\nabla\ee^{-\ii kz}\right)
 =\frac{\hbar k}{\mu}\hat{\mathbf e}_z,
\end{align*}
where $\mu$ is the reduced mass. For the scattered wave,
\[
 \mathbf j_{\rm sc}=\frac{\hbar}{2\mu\ii}
 (\psi_{\rm sc}^*\boldsymbol\nabla\psi_{\rm sc}
           -\psi_{\rm sc}\boldsymbol\nabla\psi_{\rm sc}^*).
\]
Using the spherical-coordinate gradient,
\[
 \boldsymbol\nabla
 =\hat{\mathbf r}\frac\partial{\partial r}
  +\hat{\boldsymbol\theta}\frac1r\frac\partial{\partial\theta}
  +\hat{\boldsymbol\phi}\frac1{r\sin\theta}\frac\partial{\partial\phi},
\]
we find
\[
 \boldsymbol\nabla\left(f\frac{\ee^{\ii kr}}r\right)
 =\hat{\mathbf r}\,\ii kf\frac{\ee^{\ii kr}}r+O(r^{-2}),
 \qquad
 \mathbf j_{\rm sc}\sim
 \hat{\mathbf r}\,\frac{|f|^2}{r^2}\frac{\hbar k}{\mu}.
\]
The probability flow into $\dd\Omega$ per unit time is therefore
\[
 \mathcal R(\theta,\phi;\dd\Omega)
 =\mathbf j_{\rm sc}\cdot\hat{\mathbf r}\,r^2\dd\Omega
 =|f|^2\frac{\hbar k}{\mu}\dd\Omega.
\]
Dividing by the incident flux yields
\[
 \frac{\dd\sigma}{\dd\Omega}\dd\Omega
 =\frac{\mathcal R}{j_{\rm inc}}=|f(\theta,\phi)|^2\dd\Omega,
 \qquad\boxed{\frac{\dd\sigma}{\dd\Omega}=|f(\theta,\phi)|^2.}
\]
